\pagestyle{fancy}

\section*{Introduction}
\addcontentsline{toc}{section}{Introduction}
Measuring position and motion with precision underlies a large part of modern science and technology. Semiconductor lithography, scanning probe and optical microscopy, the alignment of optical fibers and lasers, micro-electromechanical systems and biomedical devices all depend on moving objects by controlled amounts, down to micrometers and nanometers, and on knowing how far and how fast they actually moved. The devices that produce these movements, such as piezoelectric stages, flexure mechanisms and motor-driven stages, are called positioning systems, and their usefulness is determined by how well their real motion is known.

In industry and in research laboratories, the displacement of a positioning system is measured with several types of sensors. Optical incremental encoders and linear scales are widely used because they are compact and robust, although their position measurements suffer from quantization errors that make the estimation of velocity and acceleration difficult \cite{merryVelocityAccelerationEstimation2010}. Capacitive, inductive and eddy-current sensors reach a very high resolution over short ranges, and strain gauges are often built into piezoelectric stages. When the highest accuracy is required, laser interferometers are the reference: they use the wavelength of light as a ruler, which makes the measurement non-contact and traceable, and they cover long ranges with sub-nanometer resolution \cite{pedrottiIntroductionOptics2007}. Their drawbacks are cost, size and the need for a stable environment, usually an optical table with vibration isolation, so they are not always available in teaching laboratories, small research groups or developments based on low-cost hardware.

At the same time, the way in which scientific instruments are built is changing. Open-source hardware and 3D printing allow complex mechanisms to be manufactured at a low cost, and open designs such as the OpenFlexure Block Stage, a monolithic plastic flexure stage with sub-100 nm resolution, show that printed positioners can reach performances that were previously associated with machined parts \cite{mengOpenFlexureBlockStage2020, liaoLowcostOpensourceXYZ2022}. Industrial components are also being reused as scientific instruments: the optical pickup unit of a CD/DVD/Blu-ray drive integrates lasers, diffraction-limited optics and a high-bandwidth photodetector in a compact size, and has been proposed for applications such as sub-angstrom displacement sensing, imaging and force measurement \cite{hwuHackingCDDVD2018, hwuMeasurementCantileverDisplacement2006, changOpensourceForceAnalyzer2022}. On the electronics side, Field Programmable Gate Arrays (FPGAs) make it possible to acquire and time-stamp signals at hundreds of megasamples per second, which is critical in techniques such as frequency scanning interferometry \cite{sivanathanDesignFPGABasedHighSpeed2024, knollRadiationDetectionMeasurement2010}.

However, a low-cost positioner is only useful once it has been characterized. The real displacement per step of a printed stage depends on the elasticity of the plastic, the backlash and the friction of its transmission, and the way in which it is driven, and it can differ from the design value. Characterizing it requires measuring very small displacements, that is, with high spatial resolution, and the velocity with which they are reached, that is, with high temporal resolution, ideally at the same time and without relying on an expensive commercial system. The problem addressed in this work is therefore how to build an accurate, compact and low-cost instrument able to measure the displacement and the velocity of a positioning system in space and in time.

The solution proposed here combines the elements mentioned above. A compact Michelson interferometer is built from a recycled Blu-ray optical unit, so that the displacement of a movable mirror is converted into the passage of interference fringes, each one equivalent to half the wavelength of the laser. The fringe pattern is observed with two complementary detectors: a low-cost USB camera, which provides the spatial information and the direction of the motion, and a fast photodiode with an analog conditioning circuit, whose digital output is acquired and time-stamped by an FPGA and provides a high temporal resolution. All the components are held together by 3D-printed mounts, which give the system a compact and robust form that works without an optical table. The system is then used to measure the displacement and the velocity of the OpenFlexure Block Stage, a printed positioner, as a case of application.

This thesis is organized as follows. Chapter 1 presents the data acquisition system, based on an FPGA, with the communication between its processing system and its programmable logic, the ComBlock, PYNQ and the version control of the hardware design. Chapter 2 describes the Michelson interferometer: the theory of two-beam interference, its construction from the optical unit and the interference pattern obtained. Chapter 3 presents the optical detection: the camera and the method to extract the fringe motion from its images, the photodiode with its conditioning circuit, and the limits that the sampling rate imposes on the measurable velocity, together with a validation of the acquisition chain. Chapter 4 describes the experimental setup, with the 3D-printed mounts and the OpenFlexure Block Stage, and reports the measurement of its displacement and velocity.
\newpage